\emph{What is taken from the premises.} From Theorem \claimref{synthesis.null-plane-data-of-the-first-order-gravitational-field-are-the-shear-of-the-generators}, clause (iii) and the lemmas of its proof, all proved there from clauses (b), (c) and the torsion equation of the Definition for a field of any strength: (L2) for physical vector fields $A, C$ with background images $a = \hb(A)$, $c = \hb(C)$, $\hb(A \cdot D\,C - C \cdot D\,A) = (a \cdot \nabla)c - (c \cdot \nabla)a$; the derivative respects the product, $a \cdot \nabla(X \cdot Y) = (D_{a}X) \cdot Y + X \cdot (D_{a}Y)$, and is a derivation of the geometric product; (iii)(a) $\mathsf{B}$ is symmetric, $\mathsf{B}(K) = 0$, $K \cdot \mathsf{B}(A) = 0$, $K \cdot D\,K = 0$; (iii)(b) the Riccati equation $K \cdot D(\mathsf{B}(A)) - \mathsf{B}(K \cdot D\,A) = -\mathsf{B}(\mathsf{B}(A)) - \mathcal{R}(K \wedge A) \cdot K$; (iii)(c) the screen map, its trace $\theta$, its trace-free part $\sigma$, and the focusing equation. Also the form of the torsion equation used in the proof of its clause (iv): for constant background vectors $a, b$ and $g_{a} = \hb^{-1}(a)$, $D_{a}g_{b} = D_{b}g_{a}$ (this is (L2) with $A = g_{a}$, $C = g_{b}$; conversely it gives $D \wedge \hb^{T}(e^{\mu}) = 0$ by the first lemma of that proof, and then $D \wedge \hb^{T}(c) = \hb^{T}(\nabla \wedge c)$ for every field $c = c_{\mu}e^{\mu}$, so the two forms are equivalent), and the first-order curvature $-\tfrac12\,p \wedge h(B \cdot p)$ of a symmetric map $h$ at momentum $p$. From Theorem \claimref{synthesis.field-equation-linear-in-the-curvature-with-a-conserved-source-is-einsteins}, clause (vii): the form $\mathcal{G}(a) + \Lambda a = -8\pi G\,T(a)$ of the field equation, used only to read the hypotheses. For null $K$, $K \cdot \mathcal{G}(K) = K \cdot \mathcal{R}(K) - \tfrac12 Q(K)\,\mathcal{R} = K \cdot \mathcal{R}(K)$ and $K \cdot \Lambda K = 0$, so that equation gives $K \cdot \mathcal{R}(K) = -8\pi G\,K \cdot T(K)$: (E) holds when $K \cdot T(K) = 0$ and (P) when $K \cdot T(K) \geq 0$. No clause below uses the field equation otherwise.

\emph{(i).} Since $k \cdot \nabla\phi = \hb(K) \cdot \nabla\phi = K \cdot \hb^{T}(\nabla\phi) = Q(K) = 0$, the flow of $k$ preserves the hypersurfaces of constant $\phi$, and a vector carried by the flow that is tangent to one at a point stays tangent. $X_{i} \cdot K = \hb^{-1}(x_{i}) \cdot \hb^{T}(\nabla\phi) = x_{i} \cdot \nabla\phi = 0$. By (L2) with $A = K$, $C = X_{i}$, $\hb(K \cdot D\,X_{i} - X_{i} \cdot D\,K) = (k \cdot \nabla)x_{i} - (x_{i} \cdot \nabla)k = 0$, so $K \cdot D\,X_{i} = \mathsf{B}(X_{i})$ (only the derivative of $x_{i}$ along the generator and of $k$ along $x_{i}$ enter, so $x_{i}$ need be given on the generator only). $\hb^{-1}$ is an outermorphism, so $\hb^{-1}(x_{1} \wedge x_{2}) = X_{1} \wedge X_{2}$, a blade, whose square is the scalar $(X_{1} \wedge X_{2}) \cdot (X_{1} \wedge X_{2}) = (X_{1} \cdot X_{2})^{2} - Q(X_{1})Q(X_{2})$; this is minus the determinant of $G$, the matrix in the basis $X_{1}, X_{2}$ of minus the product on the screen, with entries $-X_{i} \cdot X_{j}$. The $X_{i}$ are orthogonal to the null vector $K$ and independent of it (because $x_{1} \wedge x_{2} \wedge k \neq 0$), and on the vectors orthogonal to $K$ the product is negative semidefinite with kernel spanned by $K$; hence $G$ is positive and $\rho^{4} = \det G > 0$. Adding multiples of $K$ to the $X_{i}$ changes neither $G$ nor $\rho$, so both are functions of the classes of the $X_{i}$ in the screen. The derivative is a derivation, so $K \cdot D(X_{1} \wedge X_{2}) = \mathsf{B}(X_{1}) \wedge X_{2} + X_{1} \wedge \mathsf{B}(X_{2})$, and it respects the scalar part of a product, so $(\rho^{4})' = -2\,(X_{1} \wedge X_{2}) \cdot \bigl(\mathsf{B}(X_{1}) \wedge X_{2} + X_{1} \wedge \mathsf{B}(X_{2})\bigr)$. Write $\mathsf{B}(X_{i}) = \sum_{j}\mathsf{S}_{ji}X_{j} + \alpha_{i}K$, the first term being the screen map in the basis of the classes of $X_{1}, X_{2}$. With $(a \wedge b) \cdot (c \wedge d) = (a \cdot d)(b \cdot c) - (a \cdot c)(b \cdot d)$: $(X_{1} \wedge X_{2}) \cdot (K \wedge X_{2}) = (X_{1} \cdot X_{2})(X_{2} \cdot K) - (X_{1} \cdot K)\,Q(X_{2}) = 0$ and likewise for $X_{1} \wedge K$; the terms with $X_{2} \wedge X_{2}$ and $X_{1} \wedge X_{1}$ vanish. So $(\rho^{4})' = -2(\mathsf{S}_{11} + \mathsf{S}_{22})(X_{1} \wedge X_{2})^{2} = 2\theta\rho^{4}$, the trace of a map being the same in every basis: $\rho' = \tfrac12\theta\rho$. For another pair $x_{i}' = \sum_{j}c_{ji}x_{j} + \beta_{i}k$ carried by the flow, the $c_{ji}$ are constant along the generator (apply $k \cdot \nabla$ and use that $x_{j}$, $k$ are carried and independent), and the function $\tilde\rho$ of the new pair has $\tilde\rho^{2} = |\det c|\,\rho^{2}$. For the shape: $G_{ij}' = -\mathsf{B}(X_{i}) \cdot X_{j} - X_{i} \cdot \mathsf{B}(X_{j}) = -2X_{i} \cdot \mathsf{B}(X_{j}) = 2\sum_{l}G_{il}\mathsf{S}_{lj}$, by the symmetry of $\mathsf{B}$ and $X_{i} \cdot K = 0$; so $G^{-1}G' = 2\mathsf{S} = \theta + 2\sigma$ as matrices, and with $G = \rho^{2}\bar q$, $G^{-1}G' = 2\rho'/\rho + \bar q^{-1}\bar q' = \theta + \bar q^{-1}\bar q'$: $\bar q^{-1}\bar q' = 2\sigma$ and $\ip{\sigma}{\sigma} = \tr\sigma^{2} = \tfrac14\tr\bigl((\bar q^{-1}\bar q')^{2}\bigr)$. Here $\sigma$ is the matrix, in a basis of the screen that need not be orthonormal, of the shear, a linear map of the screen; the trace of its square is that of the map. $\rho$ and $\bar q$ are defined from $\phi$ as well as from the field: a generator lies on many null hypersurfaces. The check computes $G$ from the gauge field for two separations that are not orthonormal at the point, differentiates it along the generator, and finds $\tr(G^{-1}G') = 2\theta$ and $\tfrac14\tr((G^{-1}G' - \theta)^{2}) = \ip{\sigma}{\sigma}$ on five waves that are not weak.

\emph{(ii).} From $\theta = 2\rho'/\rho$, $\theta' = 2\rho''/\rho - 2\rho'^{2}/\rho^{2} = 2\rho''/\rho - \tfrac12\theta^{2}$. The focusing equation of the cited clause (iii)(c) then reads $2\rho''/\rho = -\ip{\sigma}{\sigma} + K \cdot \mathcal{R}(K)$. For the family: with $d/d\tau = \rho^{-\alpha}\,d/d\lambda$, $d\rho^{m}/d\tau = m\rho^{m - 1 - \alpha}\rho'$ and $d^{2}\rho^{m}/d\tau^{2} = \rho^{-\alpha}\bigl(m(m - 1 - \alpha)\rho^{m - 2 - \alpha}\rho'^{2} + m\rho^{m - 1 - \alpha}\rho''\bigr)$; the term in $\rho'^{2}$ drops exactly for $m = 1 + \alpha$ (or $m = 0$), leaving $(1 + \alpha)\rho^{-\alpha}\rho'' = -\tfrac{1 + \alpha}{2}\rho^{1 - \alpha}\bigl(\ip{\sigma}{\sigma} - K \cdot \mathcal{R}(K)\bigr)$, and $\rho^{1 - \alpha}\ip{\sigma}{\sigma} = \rho^{1 + \alpha}\ip{\sigma}{\sigma}_{\tau}$. The shear taken with $\tau$ is $\rho^{-\alpha}\sigma$, since $d\bar q/d\tau = \rho^{-\alpha}\bar q'$.

\emph{(iii).} ($\alpha$) With $d/ds = \rho\,d/d\lambda$: $d\ln\rho^{2}/ds = 2\rho' = e$ and $de/ds = 2\rho\rho'' = -\rho^{2}\bigl(\ip{\sigma}{\sigma} - K \cdot \mathcal{R}(K)\bigr)$ by (ii); $\rho^{2}\ip{\sigma}{\sigma} = \ip{\sigma}{\sigma}_{s}$. Under (E), integrate once for $e$ and once more for $\ln\rho^{2}$: $\ln(\rho^{2}/\rho_{0}^{2}) = \int_{s_{0}}^{s}e = e_{0}(s - s_{0}) - \int_{s_{0}}^{s}\!\int_{s_{0}}^{t'}\ip{\sigma}{\sigma}_{s}(t)\,dt\,dt'$, and exchanging the order of integration the double integral is $\int_{s_{0}}^{s}(s - t)\ip{\sigma}{\sigma}_{s}(t)\,dt$. Then $d\lambda = \rho\,ds$. Nothing is divided by $\rho'$, so no monotonicity is needed. ($\beta$) Where $p = \rho' \neq 0$, $\rho$ is a parameter with $d/d\lambda = p\,d/d\rho$, so $\rho'' = p\,dp/d\rho = \tfrac12\,d(p^{2})/d\rho$, the shear taken with $\rho$ is $\sigma/p$, and (ii) becomes $d(p^{2})/d\rho = -\rho\,\ip{\sigma}{\sigma}_{\rho}\,p^{2} + \rho\,K \cdot \mathcal{R}(K)$, a linear equation of first order for $p^{2}$. Under (E) it is $d\ln|p|/d\rho = -\tfrac12\rho\,\ip{\sigma}{\sigma}_{\rho}$, so $p = p_{0}e^{-\Xi}$, which keeps the sign of $p_{0}$, and $d\lambda = d\rho/p$. In the ratio $v = \rho/\rho_{S}$, $\rho\,\ip{\sigma}{\sigma}_{\rho}\,d\rho = v\,\ip{\sigma}{\sigma}_{v}\,dv$, and $\ip{\sigma}{\sigma}_{v} = \tfrac14\tr\bigl((\bar q^{-1}\partial_{v}\bar q)^{2}\bigr)$ is, in components, $\tfrac14(\partial_{v}\bar q)^{AB}\partial_{v}\bar q_{AB}$ with indices raised by $\bar q$; so $\Xi = \tfrac18\int v\,(\partial_{v}\bar q)^{AB}\partial_{v}\bar q_{AB}\,dv$, the expression quoted from \cite{Chandrasekaran2026}. (A labelled translation into components, for the comparison only.) In both forms the datum is the shape as a function of a parameter built from the solution ($s$ is defined through $\rho$, and in ($\beta$) $\rho$ is the parameter); if the shape is given as a function of $\lambda$ the formulas are integral equations for $\rho$, equivalent to the linear equation of (ii), and $\lambda$ follows by the last quadrature. The check verifies both forms, with and without the term for matter, as identities between power series along a generator of exact plane waves, to order eight, with the series of $\theta$, $\ip{\sigma}{\sigma}$ and $\rho$ computed from the gauge fields; form ($\beta$) is verified with $\rho$ itself as the variable, $\lambda$ being reverted as a series in $\rho$, including the quadrature for $\lambda$.

\emph{(iv).} By (ii) and (P), $\rho'' = -\tfrac12\bigl(\ip{\sigma}{\sigma} - K \cdot \mathcal{R}(K)\bigr)\rho \leq 0$ where $\rho > 0$, both terms of the bracket being non-negative. A concave function lies below its tangents: $\rho(\lambda) \leq \rho_{0} + \rho'(\lambda_{0})(\lambda - \lambda_{0}) = \rho_{0}\bigl(1 + \tfrac12\theta_{0}(\lambda - \lambda_{0})\bigr)$. If the generator stayed where the field and $\phi$ are smooth, with $\rho$ smooth and positive, on all of $[\lambda_{0}, \lambda_{0} + 2/|\theta_{0}|]$ with $\theta_{0} < 0$, the right side would vanish at the end point, a contradiction. So that stretch ends within the interval. The argument does not show how it ends: that $\phi$, and with it $K$ and $\theta$, stay smooth as long as the field is smooth and $\rho > 0$ would need the equation of the separations themselves, which is used in (vii) for the plane waves and is not proved here for a general field. Equality at some $\lambda > \lambda_{0}$ forces $\rho'' = 0$ on $[\lambda_{0}, \lambda]$, that is, both non-negative terms zero there. $\rho'$ is non-increasing, and $\theta = 2\rho'/\rho$ has its sign. If $\rho$ is constant on an interval then $\rho'' = 0$ there and again both terms vanish.

\emph{(v).} Given $(\bar q(s), \rho_{0}, e_{0})$ on an interval $I$ of $s$ containing $s_{0}$: $\ip{\sigma}{\sigma}_{s} = \tfrac14\tr((\bar q^{-1}d\bar q/ds)^{2})$ is a smooth function on $I$, the formulas of (iii)($\alpha$) define $e$ and a positive $\rho$ on all of $I$, and $\lambda(s)$ is strictly increasing, so $s$ is a smooth function of $\lambda$ and $(\rho, \bar q)$ become functions of $\lambda$ on an interval; they satisfy $2\rho\rho'' = de/ds = -\rho^{2}\ip{\sigma}{\sigma}$, the equation of (ii) under (E), with $\ip{\sigma}{\sigma} = \rho^{-2}\ip{\sigma}{\sigma}_{s} = \tfrac14\tr((\bar q^{-1}\bar q')^{2})$. Conversely, given $(\rho(\lambda), \bar q(\lambda))$ with $\rho > 0$ satisfying that equation, $s = s_{0} + \int d\lambda/\rho$ is strictly increasing, $\bar q$ becomes a function of $s$, and $\rho_{0} = \rho(\lambda_{0})$, $e_{0} = 2\rho'(\lambda_{0})$; by (iii)($\alpha$) the formulas return $\rho$. The two constructions are inverse to each other. On a stretch where $\rho$ is monotonic the same holds with $\bar q(\rho)$ and (iii)($\beta$), the constants being $\rho_{0}$ and $p_{0} \neq 0$ and the interval one on which $\bar q(\rho)$ is given. The last sentence of the clause states what is not proved.

\emph{(vi)(a).} $\ell$ is null and $W$ depends on $x$ through $\lambda = \ell \cdot x$ and $y$. $\hb(\hb^{-1}(a)) = a - \tfrac12 W(a \cdot \ell)\ell + \tfrac12 W(a \cdot \ell)\ell = a$, because $\ell \cdot \ell = 0$; in a basis containing $\ell$ the matrix of $\hb$ is the identity plus a multiple of one off-diagonal entry, so $\det\hb = 1$. $\hb$ is symmetric, $\hb(\ell) = \ell$ and $\hb(a) \cdot \ell = a \cdot \ell$. For a vector $a$ write $a_{\perp} = a - (a \cdot \ell)n - (a \cdot n)\ell$; then $(a \cdot \nabla)y = a_{\perp}$, $(a \cdot \nabla)\lambda = a \cdot \ell$, and, $\mathsf{w}$ being symmetric, $\nabla W = -2\,\mathsf{w}(y) - (y \cdot \mathsf{w}'(y))\,\ell$. With $g_{b} = b - \tfrac12 W(b \cdot \ell)\ell$ and $(u \wedge v) \cdot c = u\,(v \cdot c) - v\,(u \cdot c)$:
\[
(a \cdot \nabla)g_{b} = \bigl((a \cdot \mathsf{w}(y)) + \tfrac12(a \cdot \ell)\,y \cdot \mathsf{w}'(y)\bigr)(b \cdot \ell)\,\ell, \qquad \Omega(a) \cdot g_{b} = (a \cdot \ell)\bigl((b \cdot \mathsf{w}(y))\,\ell - (b \cdot \ell)\,\mathsf{w}(y)\bigr),
\]
the second because $(\ell \wedge \mathsf{w}(y)) \cdot \ell = 0$. Their sum $D_{a}g_{b}$ is symmetric in $a, b$, which is the torsion equation in the form recalled above; it fixes $\Omega$ (Definition, clause (f)). For the field strength of clause (c) of the Definition, with constant $a, b$: $\Omega(a) \times \Omega(b)$ is a multiple of $(\ell \wedge \mathsf{w}(y)) \times (\ell \wedge \mathsf{w}(y)) = 0$, and $(a \cdot \nabla)\Omega(b) = (b \cdot \ell)\;\ell \wedge \bigl(\mathsf{w}(a_{\perp}) + (a \cdot \ell)\mathsf{w}'(y)\bigr)$; in the difference the terms with $\mathsf{w}'$ cancel and $R(a \wedge b) = \ell \wedge \mathsf{w}\bigl((b \cdot \ell)a - (a \cdot \ell)b\bigr) = \ell \wedge \mathsf{w}((a \wedge b) \cdot \ell)$, with $\mathsf{w}$ extended by zero on $n$ and $\ell$. Since $\hb(a \wedge b) \cdot \ell = \hb(a)(b \cdot \ell) - \hb(b)(a \cdot \ell) = (a \wedge b) \cdot \ell$ (the two terms in $\ell$ cancel), $\mathcal{R}(B) = -R(\hb(B)) = -\ell \wedge \mathsf{w}(B \cdot \ell)$, which is $-\tfrac12\,p \wedge h(B \cdot p)$ at $p = \ell$, $h = 2\mathsf{w}$; $h$ is transverse and symmetric at the null momentum $\ell$, trace-free exactly when $\tr\mathsf{w} = 0$. Contracting, $\mathcal{R}(b) = -e^{\mu} \cdot \bigl(\ell \wedge \mathsf{w}(e_{\mu}(b \cdot \ell) - b\,(e_{\mu} \cdot \ell))\bigr) = -(b \cdot \ell)\bigl(\mathsf{w}(\ell) - \ell\,(e^{\mu} \cdot \mathsf{w}(e_{\mu}))\bigr) + \ell \cdot (\ell \wedge \mathsf{w}(b)) = (\tr\mathsf{w})(b \cdot \ell)\,\ell$, using $e^{\mu}(e_{\mu} \cdot \ell) = \ell$, $\mathsf{w}(\ell) = 0$ and $\ell \cdot (\ell \wedge v) = Q(\ell)v - (\ell \cdot v)\ell = 0$ for transverse $v$. With $\tr\mathsf{w} = 0$ the Ricci map, hence the Einstein map, vanishes. The derived metric is $g_{a} \cdot g_{b} = a \cdot b - W(a \cdot \ell)(b \cdot \ell)$. The check solves $\Omega$ from the torsion equation with the gauge fields as jets, on five profiles that are not small at a point that is not the origin, and finds the closed form, the curvature on all six planes and the Ricci map as stated.

\emph{(vi)(b).} $\nabla\phi = n + \mathsf{S}(y) + \tfrac12(y \cdot \mathsf{S}'(y))\,\ell$, and $\nabla\phi \cdot \ell = 1$, so $K = \hb^{T}(\nabla\phi) = \nabla\phi + \tfrac12 W\ell$, as stated, and $Q(K) = Q(\mathsf{S}(y)) + \bigl(y \cdot \mathsf{S}'(y) + W\bigr)(n \cdot \ell) = y \cdot (\mathsf{S}^{2} + \mathsf{S}' - \mathsf{w})(y)$, which vanishes for every $y$ exactly when $\mathsf{S}' = \mathsf{w} - \mathsf{S}^{2}$, a symmetric map being fixed by its quadratic form. $k = \hb(K) = K + \tfrac12 W\ell$, so $\ell \cdot k = K \cdot \ell = 1$: $d\lambda/d(\text{affine parameter}) = 1$, and the transverse part of $k$ is $\mathsf{S}(y)$, so along a generator $y' = \mathsf{S}(y)$ and $y = E(y_{0})$. For transverse $e$: $K \cdot m(e) = K \cdot e - (K \cdot e)(K \cdot \ell) = 0$; $m(e) \cdot m(e') = e \cdot e'$ since $\ell$ is null and orthogonal to $e, e'$; and $m(e) \cdot \ell = 0$, so $\hb(m(e)) = m(e)$. Next $K \cdot D\,m(e) = (k \cdot \nabla)m(e) + \Omega(k) \cdot m(e)$; the first term is a multiple of $\ell$, and $\Omega(k) \cdot m(e) = (\ell \wedge \mathsf{w}(y)) \cdot m(e) = \ell\,(\mathsf{w}(y) \cdot e)$, also a multiple of $\ell$; but $K \cdot (K \cdot D\,m(e)) = k \cdot \nabla(K \cdot m(e)) - (K \cdot D\,K) \cdot m(e) = 0$ while $K \cdot \ell = 1$, so the multiple is zero: $K \cdot D\,m(e) = 0$. For $\mathsf{B}$: $\Omega(m(e)) = 0$ because $m(e) \cdot \ell = 0$, and $(m(e) \cdot \nabla)K = \mathsf{S}(e) + c\,\ell$ for a scalar $c$ (only the transverse part $e$ of $m(e)$ moves $y$, and $m(e) \cdot \nabla\lambda = 0$); $\mathsf{B}(m(e))$ is orthogonal to $K$ by (iii)(a), so $c = -K \cdot \mathsf{S}(e)$ and $\mathsf{B}(m(e)) = m(\mathsf{S}(e))$. The classes of $m(e_{1}), m(e_{2})$ span the screen, so the screen map is $\mathsf{S}$ in that orthonormal frame, with the stated $\theta$ and $\sigma$. From (a), $K \cdot \mathcal{R}(K) = (\tr\mathsf{w})(K \cdot \ell)^{2} = \tr\mathsf{w}$; $(K \wedge m(e)) \cdot \ell = -m(e)$, so $\mathcal{R}(K \wedge m(e)) = \ell \wedge \mathsf{w}(e)$, $(K \wedge m(e')) \cdot (\ell \wedge \mathsf{w}(e)) = (K \cdot \mathsf{w}(e))(m(e') \cdot \ell) - (K \cdot \ell)(m(e') \cdot \mathsf{w}(e)) = -e' \cdot \mathsf{w}(e)$, and $\mathcal{R}(K \wedge m(e)) \cdot K = \ell\,(\mathsf{w}(e) \cdot K) - \mathsf{w}(e) = -m(\mathsf{w}(e))$. The Riccati equation of (iii)(b) at the parallel field $A = m(e)$ is then $m(\mathsf{S}'(e)) = -m(\mathsf{S}^{2}(e)) + m(\mathsf{w}(e))$: the same equation as the condition that $K$ be null. The background vectors $x_{i} = m(E(e_{i}))$ are tangent ($x_{i} \cdot \nabla\phi = K \cdot m(E(e_{i})) = 0$, as $\hb^{-1}(x_{i}) = x_{i}$) and are the derivatives of the position on a generator with respect to $y_{0}$ at fixed $\lambda$ and $\phi$, hence carried by the flow; $-x_{i} \cdot x_{j} = -E(e_{i}) \cdot E(e_{j})$ is the matrix $E^{T}E$, so $\rho^{4} = (\det E)^{2}$. If $\mathsf{S}(\lambda_{0}) = 0$ then $\theta$ and $\sigma$ vanish at $\lambda_{0}$ and $K = n + \tfrac12(y \cdot \mathsf{w}(y) + W)\ell = n$ there; and if $\mathsf{w} = 0$ on an interval ending at $\lambda_{0}$, $\mathsf{S}' = -\mathsf{S}^{2}$ with $\mathsf{S}(\lambda_{0}) = 0$ gives $\mathsf{S} = 0$ on it, so $\phi = n \cdot x$ there. The check verifies each statement of (b) on four waves that are not weak, as jets, with $\Omega$ solved from the torsion equation, and (i) and (ii) on five.

\emph{(vi)(c).} In two dimensions a trace-free symmetric map has $\hat\sigma^{2} = \tfrac12\tr(\hat\sigma^{2})\,\id$. So $\mathsf{S}^{2} = \tfrac14\theta^{2} + \theta\hat\sigma + \tfrac12\tr\hat\sigma^{2}$, and $\mathsf{w} - \mathsf{S}^{2} = \hat\sigma' - \tfrac14\theta^{2} - \tfrac12\tr\hat\sigma^{2} = \hat\sigma' + \tfrac12\theta' = \mathsf{S}'$ by the equation for $\theta$; $\tr\mathsf{w} = \tr\hat\sigma' + \theta\tr\hat\sigma = 0$. By (b) the hypersurface has shear $\hat\sigma$ and expansion $\theta$ in the frame $m(e_{1}), m(e_{2})$. For the second statement let $G = \rho^{2}\bar q$ and $A = G^{1/2}$, the positive symmetric root, smooth in $\lambda$. The antisymmetric part of $A'A^{-1}$ is $a\,J$ for a function $a$, $J$ the quarter turn of the plane (the action of the bivector $e_{1}e_{2}$ on its vectors). Let $O$ be the rotation with $O' = -a\,J\,O$, and $E = OA$. Then $E^{T}E = G$ and $\mathsf{S} = E'E^{-1} = -a\,J + O\,A'A^{-1}\,O^{T}$, whose antisymmetric part is $-aJ + aJ = 0$ since a rotation commutes with $J$: $\mathsf{S}$ is symmetric, and $\mathsf{w} = \mathsf{S}' + \mathsf{S}^{2} = E''E^{-1}$ is symmetric. $\tr\mathsf{S} = \tr(A'A^{-1}) = (\ln\det A)' = 2\rho'/\rho$, and by (i) applied to the separations $m(E(e_{i}))$ (see below) the square of the shear is $\tfrac14\tr((\bar q^{-1}\bar q')^{2})$; so $\tr\mathsf{w} = \theta' + \tfrac12\theta^{2} + \ip{\sigma}{\sigma} = (2/\rho)\bigl(\rho'' + \tfrac12\ip{\sigma}{\sigma}\rho\bigr) = 0$. Here $E(\lambda_{0}) = O(\lambda_{0})A(\lambda_{0})$ is not the identity, so (b) is applied as follows. With $E_{b}$ the map of (b) ($E_{b}(\lambda_{0}) = \id$) for this $\mathsf{S}$, $E = E_{b}C$ with $C = E(\lambda_{0})$ constant, both sides solving $E' = \mathsf{S}E$ with the same value at $\lambda_{0}$; so the vectors $x_{i} = m(E(e_{i}))$ are linear combinations with constant coefficients of the separations of (b), hence carried separations, with Gram matrix $E^{T}E = G$, and (i) applies to them. So the hypersurface with this $\mathsf{S}$ in the wave with this $\mathsf{w}$ has $G = E^{T}E$ for these separations. \emph{The orders.} With $\mathsf{S}(\lambda_{0}) = 0$ and $\tr\mathsf{w} = 0$, $\mathsf{S} = \int\mathsf{w} + O(\mathsf{w}^{2})$ is trace-free at first order, so $\sigma = \int\mathsf{w}$, $\theta = 0$ at that order, and the trace of the Riccati equation gives $\theta' = -\ip{\sigma}{\sigma}$ at second order. Comparison with clause (iv) of the first-order theorem: there $\sigma = \tfrac12\partial_{n}h^{TT}$ and the curvature fixes $\partial_{n}\sigma$; here $\partial_{\lambda}\sigma = \mathsf{w}$ at first order with curvature that of $h = 2\mathsf{w}$ at $p = \ell$, for which $\partial_{n}^{2}h = (n \cdot \ell)^{2}\,2\mathsf{w}$: the same. The check takes $\mathsf{S} = 0$ at the point, six amplitudes of one profile with both polarisations, and jets of order eight; each coefficient of the series of $\theta$ and $\rho$ in $\lambda$ through that order is a polynomial of degree at most four in the amplitude, recovered by interpolation, the sixth amplitude confirming the degree. In vacuum the first-order coefficient of $\rho$ vanishes, the second-order one is $-\tfrac{61}{48}(\lambda - \lambda_{0})^{4}$ at lowest order (minus $\tr\mathsf{w}^{2}/24$ with $\tr\mathsf{w}^{2} = \tfrac{61}{2}$ at the point), and the odd orders vanish identically because the quarter turn $J$ carries $\mathsf{w}$ to $-\mathsf{w}$ and leaves $\theta$ and $\rho$ unchanged; on a wave with matter all four orders are present. The equations of (ii) and (iii) hold on every series. A further check takes exact waves $\mathsf{w} = \mathrm{diag}(c_{1}, c_{2})/\lambda^{2}$: for the shape $\bar q = \mathrm{diag}(\rho^{a}, \rho^{-a})$ form ($\beta$) gives $\Xi = \tfrac{a^{2}}{4}\ln(\rho/\rho_{0})$, $p \propto \rho^{-a^{2}/4}$ and $\lambda \propto \rho^{1 + a^{2}/4}$, so $E = \mathrm{diag}(\lambda^{p_{1}}, \lambda^{p_{2}})$ with $p_{1,2} = (1 \pm \tfrac{a}{2})/(1 + \tfrac{a^{2}}{4})$ and $c_{i} = p_{i}(p_{i} - 1)$, $c_{1} + c_{2} = 0$. The check tests this in differential form at two sample points for six values of $a$: from the gauge fields of that wave, the Gram matrix of the separations has $\bar q^{-1}\bar q' = \tfrac{a\theta}{2}\,\mathrm{diag}(1, -1)$, and $d\ln p/d\ln\rho = -\tfrac{a^{2}}{4}$, $d\ln\lambda/d\ln\rho = 1 + \tfrac{a^{2}}{4}$ there.

\emph{(vii).} $E$ solves $E'' = (\mathsf{S}' + \mathsf{S}^{2})E = \mathsf{w}E$, a linear equation with smooth coefficients, so it exists on all of $[\lambda_{0}, \lambda_{2}]$, and $\mathsf{S} = E'E^{-1}$, the hypersurface and $\rho^{2} = |\det E|$ are smooth exactly while $\det E \neq 0$. Suppose $\det E \neq 0$ on $[\lambda_{0}, \lambda_{1}]$. The trace of the Riccati equation is $\theta' = -\tfrac12\theta^{2} - \ip{\sigma}{\sigma} + \tr\mathsf{w} \leq -(\ip{\sigma}{\sigma} - \tr\mathsf{w})$ with $\theta(\lambda_{0}) = 0$, so $\theta(\lambda_{1}) \leq -\int(\ip{\sigma}{\sigma} - \tr\mathsf{w})$, and the integrand is non-negative. If it vanished identically on $[\lambda_{0}, \lambda_{1}]$ then $\sigma = 0$ and $\tr\mathsf{w} = 0$ there, $\mathsf{S} = \tfrac12\theta\,\id$, and the trace-free part of $\mathsf{S}' = \mathsf{w} - \mathsf{S}^{2}$ would give $\mathsf{w} = 0$ on the interval, against the hypothesis. So $\theta(\lambda_{1}) < 0$. The wave obeys (P) for this $K$ since $K \cdot \mathcal{R}(K) = \tr\mathsf{w} \leq 0$, and by (iv) the stretch on which the hypersurface is smooth with $\rho > 0$ ends by $\lambda_{1} + 2/|\theta(\lambda_{1})|$; if that value is at most $\lambda_{2}$, the field and $E$ are smooth up to it, the hypersurface is smooth while $\det E \neq 0$, and so the stretch can end only by $\det E = 0$: $\rho$ vanishes there. There $\hb$ is smooth with $\det\hb = 1$ and the Ricci map is $(\tr\mathsf{w})(a \cdot \ell)\ell$. This is the focusing property of plane waves of \cite{Penrose1965}. The check exhibits exact cases. (1) A wave with $\tr\mathsf{w} < 0$ crossing a null plane with $\rho^{2} = 1 - \tfrac32\lambda^{2} + \tfrac12\lambda^{3}$, zero at $\lambda = 1$ within the bound from each of four points. (2) Waves without matter, smooth for every $\lambda$, with a column of $E$ vanishing at $\lambda = 1$ where the field is smooth, on hypersurfaces that are not planes at the start. (3) A wave without matter taking a null plane to a caustic: $\mathsf{w} = \mathrm{diag}(6/q, -6/q)$ with $q = 2 + 2\lambda - \lambda^{2}$, smooth and positive on $[0, 1]$, and $\mathsf{S}(0) = 0$. Then $E = \mathrm{diag}(f, b)$ with $b = 1 - \tfrac32\lambda^{2} + \tfrac12\lambda^{3}$ ($b'' = -(6/q)\,b$, $b(0) = 1$, $b'(0) = 0$, $b(1) = 0$, $b'(1) = -\tfrac32$) and $f'' = (6/q)f$, $f(0) = 1$, $f'(0) = 0$; since $6/q > 0$ on $[0, 1]$, $f'' > 0$ while $f > 0$, so $f' \geq 0$ and $f \geq 1$ there, and $\det E = f\,b$ is positive on $[0, 1)$ and vanishes at $\lambda = 1$, where the field is smooth, invertible and without matter. The check verifies the series of $E$ to order eight (the column $b$ exactly, the off-diagonal entries zero), the field and the null plane at $\lambda = 0$ from the gauge fields, and the field at $\lambda = 1$. (4) The power-law waves, on which the bound is attained exactly when the shear vanishes.

\emph{(viii).} Replacing $\phi$ by $\phi/a$ replaces $K$ by $K/a$ and $\lambda$ by $a\lambda$, and the origin of $\lambda$ on a generator is free; by (i) $\rho$ is fixed up to a constant factor. Under both, $p_{0} = \rho'(\lambda_{0})$ becomes $(c/a)p_{0}$, so $v$ and $t$ are unchanged. From (iii)($\beta$), $\lambda - \lambda_{0} = (\rho_{0}/p_{0})\int_{1}^{v}e^{\Xi}\,dr$ and $\Xi = \tfrac12\int_{1}^{v}r\,\ip{\sigma}{\sigma}_{r}\,dr$ in the ratio, which is the stated $t(v)$; the integrand of $\Xi$ is non-negative, so $\Xi \geq 0$ for $v \geq 1$ and $\Xi \leq 0$ for $v \leq 1$, giving the two inequalities. On a growing stretch $p_{0} > 0$ and $\lambda - \lambda_{0} = \rho_{0}t/p_{0} \geq \rho_{0}(v - 1)/p_{0}$; on a falling one $p_{0} < 0$, $v < 1$ and $\lambda - \lambda_{0} = \rho_{0}|t|/|p_{0}| \leq \rho_{0}(1 - v)/|p_{0}|$. The statements under (P) are those of (iv): $\rho'$ non-increasing gives at most one sign change of $p$, from positive to negative, and constancy of $\rho$ on an interval gives $\sigma = 0$, $K \cdot \mathcal{R}(K) = 0$ there. For $s$: $ds/d\lambda = 1/\rho > 0$, and $de/ds = -\ip{\sigma}{\sigma}_{s} + \rho^{2}K \cdot \mathcal{R}(K) \leq 0$ under (P). The check verifies the relation of $\lambda$ to $\rho$ on the power-law waves, where $d\ln\lambda/d\ln\rho = 1 + \tfrac{a^{2}}{4}$.

\emph{Reading.} With the field equation, $K \cdot \mathcal{R}(K) = -8\pi G\,K \cdot T(K)$, and clause (ii) is $\rho'' = -\tfrac12\bigl(\ip{\sigma}{\sigma} + 8\pi G\,K \cdot T(K)\bigr)\rho = -4\pi G\bigl(\ip{\sigma}{\sigma}/8\pi G + K \cdot T(K)\bigr)\rho$; the identification of $\ip{\sigma}{\sigma}/8\pi G$ with a density of null energy is clause (vi) of the first-order theorem and holds at the leading order.
